En 1948, John von Neumann, con ayuda de Stanislaw Ulam, propuso un autómata celular como un modelo matemático con la capacidad de autorreproducción universal, el cual desarrolló en su libro \textit{Theory of Self-Reproducing Automata} \cite{vonNeumann1948}. Hoy en día representa un modelo fundamental para el estudio de sistemas complejos, computación distribuida y procesos autoorganizados. 

Varios autores, como Stephen Wolfram \cite{wolfram1982} y Erick W. Weisstein \cite{weissteinCA}, están de acuerdo en describir los autómatas celulares como sistemas discretos formados por células que se organizan en una retícula\footnote{Conjunto de líneas o elementos de estructura lineal que se cruzan o entrelazan estrechamente \cite{fundacionBBVALattice}} (\textit{lattice}). Cada celda puede tener un valor de un número limitado de estados, y se actualiza en intervalos de tiempo específicos. Esta actualización se basa en una función de transición que solo depende de las celdas cercanas.

Aunque existen variaciones en la terminología (celda, célula, sitio, color), la estructura formal es ampliamente aceptada y utilizada en campos como la física, biología computacional, procesamiento de imágenes, teoría de la computación, entre otros. Más recientemente, se han integrado técnicas de aprendizaje supervisado, en particular redes neuronales artificiales, dando lugar a los autómatas celulares neuronales.

Este estado del arte revisa estas propuestas combinadas, destacando su uso en la creación de formas artificiales, la clasificación de patrones, la generación de texturas en movimiento y el modelado en 3D. Además, se encuentran desafíos relacionados con la capacidad de escalar, la posibilidad de aplicar lo aprendido en diferentes situaciones y el control de comportamientos que surgen de manera inesperada. Estos aspectos impulsan la creación de nuevas versiones de los autómatas celulares neuronales.